\documentclass{article}

\usepackage{amsmath,amssymb}
\usepackage{xurl}
\usepackage[colorlinks=true,linkcolor=blue,citecolor=blue,urlcolor=blue]{hyperref}
\setlength{\emergencystretch}{2em}

% Shared commands for notes in docs/paper-gaps/.

\DeclareUnicodeCharacter{2080}{\ifmmode{}_{0}\else\textsubscript{0}\fi}
\DeclareUnicodeCharacter{2081}{\ifmmode{}_{1}\else\textsubscript{1}\fi}
\DeclareUnicodeCharacter{2082}{\ifmmode{}_{2}\else\textsubscript{2}\fi}
\DeclareUnicodeCharacter{2099}{\ifmmode{}_{n}\else\textsubscript{n}\fi}

\newcommand{\Lean}{\textsc{Lean}}
\ExplSyntaxOn
\NewDocumentCommand{\leanid}{m}{
  \ifmmode
    \textsf{\detokenize{#1}}
  \else
    \group_begin:
    \urlstyle{sf}
    \tl_set:Nn \l_tmpa_tl {#1}
    \tl_replace_all:Nnn \l_tmpa_tl {\_} {_}
    \exp_args:NV \path \l_tmpa_tl
    \group_end:
  \fi
}
\ExplSyntaxOff
\providecommand{\tr}{\operatorname{tr}}
\newcommand{\tnleanIssueBase}{https://github.com/LionSR/TNLean/issues/}
\newcommand{\tnleanPrBase}{https://github.com/LionSR/TNLean/pull/}
\newcommand{\tnleanDocBase}{https://sirui-lu.com/TNLean/docs/}
\newcommand{\ghissue}[1]{\href{\tnleanIssueBase#1}{GitHub issue \##1}}
\newcommand{\ghpr}[1]{\href{\tnleanPrBase#1}{PR \##1}}
\newcommand{\doclink}[1]{\href{\tnleanDocBase#1.html}{\path{#1}}}

% Dirac notation and the identity operator, as in blueprint/src/macros/common.tex.
% \providecommand keeps note-local definitions authoritative.
\usepackage{dsfont}
\providecommand{\ket}[1]{|#1\rangle}
\providecommand{\bra}[1]{\langle #1|}
\providecommand{\braket}[2]{\langle #1 | #2 \rangle}
\providecommand{\ketbra}[2]{|#1\rangle\!\langle #2|}
\providecommand{\Id}{\mathds{1}}
\providecommand{\one}{\mathds{1}}
\providecommand{\C}{\mathbb{C}}
\providecommand{\MN}[1]{M_{#1}(\C)}
\providecommand{\MD}{M_D(\C)}

% External referencing for paper sources.  The build helper writes sanitized
% auxiliary files under build/paper-aux/ and excludes duplicated source labels.
\usepackage{xr}

% Import a generated paper auxiliary file with a caller-supplied label prefix.
% The candidate paths support builds from docs/paper-gaps/, the repository root,
% and build/paper-aux/ itself.
\newcommand{\importpaperaux}[2]{%
  \IfFileExists{../../build/paper-aux/#2.aux}{%
    \externaldocument[#1]{../../build/paper-aux/#2}%
  }{%
    \IfFileExists{build/paper-aux/#2.aux}{%
      \externaldocument[#1]{build/paper-aux/#2}%
    }{%
      \IfFileExists{#2.aux}{%
        \externaldocument[#1]{#2}%
      }{}%
    }%
  }%
}

% General external reference macros
\newcommand{\paperref}[2]{\ref{#1-#2}}
\newcommand{\papereqref}[2]{(\ref{#1-#2})}

% CPSV16 source (1606.00608 / MPDO-22-12-17-2.tex) external document and shortcuts
\importpaperaux{cpsv16-}{1606.00608/MPDO-22-12-17-2-tnlean}
\newcommand{\cpsvref}[1]{\paperref{cpsv16}{#1}}
\newcommand{\cpsveqref}[1]{\papereqref{cpsv16}{#1}}

% Verdict marker: \gapnote{<kind>}{<status>}. Kind is the policy
% classification (clarification, local-correction, scope-restriction,
% unfaithful, false-source, open-gap); status is open, wip (elimination
% underway), resolved, or historical. Machine-read by the site index;
% prints a verdict line.
\newcommand{\gapnote}[2]{\par\noindent\textbf{Verdict:} #1 (#2).\par}


% Fallbacks keep this author document compilable when it is built outside its
% source directory and a different command.tex is found first.
\providecommand{\Lean}{\textsc{Lean}}
\providecommand{\leanid}[1]{\textsf{\detokenize{#1}}}
\providecommand{\doclink}[1]{%
    \href{https://sirui-lu.com/TNLean/docs/#1.html}{\path{#1}}}
\providecommand{\gapnote}[2]{%
    \par\noindent\textbf{Verdict:} #1 (#2).\par}
\providecommand{\tr}{\operatorname{tr}}
\providecommand{\braket}[2]{\langle #1\mid #2\rangle}

\title{Policy for Paper-Gap Notes}
\author{}
\date{2026-05-12}

\begin{document}
\maketitle

\section{Purpose}

A paper-gap note records a mathematical discrepancy between a cited source and
the statement, proof, or route used in the formal development. The archived
source itself must remain intact. When a printed statement lacks the stable
anchor needed for an external reference, a build helper may add a metadata-only
\verb|\label| to a temporary copy under \path{build/}; never add the label to
the archived source. If a proof appears to require a hypothesis absent from the
article, a different scalar, a different normalization, or a replacement
theorem from the formal library, explain the change in a separate note under
\path{docs/paper-gaps/}.

The note is a mathematical document first. It should be readable by a
mathematician or mathematical physicist who has not read the corresponding
issue or pull request. Traceability data such as formal declaration names,
source-file paths, GitHub issues, and line ranges support the exposition,
usually in footnotes; they do not govern it.

\section{When a note is needed}

A note is needed when the formal statement is not literally the cited
statement. A common case is that the paper claims
\begin{align}
    \mathcal H &\Longrightarrow C.
    \label{eq:policy_source_implication}
\end{align}
The formal proof may instead establish only
\begin{align}
    \mathcal H\wedge\mathcal K &\Longrightarrow C.
    \label{eq:policy_formal_implication}
\end{align}
The condition $\mathcal K$ must then be accounted for. Either it follows from
$\mathcal H$, it is a harmless local correction, or
(\ref{eq:policy_formal_implication}) is a restricted theorem relative to
(\ref{eq:policy_source_implication}) and must be marked as such.

A note is also needed when the proof is mathematically different even though
the final conclusion is the same. For example, a source proof might identify a
coefficient from the following finite power-sum identity, assumed for
$1\leq N\leq m$:
\begin{align}
    \sum_{q=1}^{r_a}\mu_q^N
        &= \sum_{q=1}^{r_b}\nu_q^N.
    \label{eq:policy_power_sum}
\end{align}
The formal route might instead use convergence of a sequence. These are
different mechanisms. The note must state which hypotheses each mechanism uses
and show the intermediate equalities required by the replacement argument.

A note is \emph{not} an invitation to model a degenerate reading. When a source
definition, read literally, admits a case its authors do not intend, such as a
direct-sum summand with coefficient zero, a basis member whose coefficient
vanishes at every length, a sector of probability zero, or an empty block
padding a bond dimension, adopt the intended reading as a convention and write
it into the definition. One short note of kind \path{local-correction} records
that reading. The degenerate case receives no parallel definition,
counterexample, status entry in the blueprint, or hypothesis repeated on
downstream statements; delete such apparatus when it is found. Reserve a note
of kind \path{false-source} for a printed claim that fails on nondegenerate
data.

\section{Mathematical presentation}

Define every symbol before its first use. State the cited assertion in the
special case under discussion, state the formal or corrected assertion, and
compare their hypotheses and conclusions. Use precise, concise, and
mathematically direct prose. Long quotations are usually unnecessary; a
specific citation and the relevant formulas are more useful.

Every substantive displayed equation must use
\verb|\begin{align}| and \verb|\end{align}|. Use multiple aligned lines when a
derivation depends on intermediate equalities. An \verb|aligned| structure may
occur inside \verb|align| only when a nested multi-line expression is
mathematically necessary. Do not compress a proof gap into one implication when
the missing step occurs within a calculation.

Every substantive displayed equation must have a stable semantic label of the
form \verb|\label{eq:<note-prefix>_<content>}|. The label must describe the
mathematical content, not the display's position or current number. When a
display contains several independently discussed lines, label each such line.
For example, a comparison argument may use
\begin{align}
    V^{(N)}(A)
        &= c_N V^{(N)}(B),
    \label{eq:policy_proportional_vectors}\\
    \braket{V^{(N)}(B_k)}{V^{(N)}(A_j)}
        &\xrightarrow[N\to\infty]{}0,
    \label{eq:policy_mixed_overlap}\\
    \dim\operatorname{span}
        \left(\{V^{(N)}(A_j)\}_j\cup\{V^{(N)}(B_k)\}\right)
        &= r_A+1.
    \label{eq:policy_span_dimension}
\end{align}
The prose must identify the nonzero coefficient forced to vanish and cite the
step that produces the contradiction, such as
(\ref{eq:policy_mixed_overlap}) or
(\ref{eq:policy_span_dimension}).

Refer to a displayed equation in the local note through \verb|\ref|. Write
\verb|(\ref{eq:<label>})| or a grammatical equivalent. Do not write ``the
equation above'', ``the preceding display'', a raw local equation number, or an
unlinked source equation number. A source-paper equation is not a local
cross-reference: name the source, identify its equation, and cite it with an
existing bibliography key, for example Equation~(5.2) of
Ref.~\cite{Wolf2012Quantum}. When a local \LaTeX{} source is available, run the
paper's declared auxiliary-label build step before compiling its paper-gap
notes; for CPSV16, run
\path{./scripts/generate_paper_aux.sh 1606.00608}. Cite source equations,
theorems, and definitions through a prefixed external label so the rendered
number comes from the source; do not hardcode numbers or expose raw source label
names. The build helper imports only labels that occur exactly once and excludes
ambiguous duplicates. A note may compile without the generated auxiliary file,
but unresolved external references render as ``??'' and are not acceptable in a
published artifact. If the source statement has no stable label or auxiliary
file, use an ordinary bibliographic citation and explicit printed designation
rather than inventing a number. Every mathematical
claim attributed to a paper or other literature source must likewise carry a
\verb|\cite{...}| citation using a key already present in \path{references.bib}.
Do not invent citation keys.

Do not use \verb|$$|\ldots\verb|$$|, \verb|\[|\ldots\verb|\]|,
\verb|equation|, \verb|equation*|, \verb|align*|, \verb|gather|,
\verb|gather*|, \verb|multline|, or \verb|multline*| for equations in
paper-gap notes. In particular, an equation later discussed in prose may not
appear in an unlabelled display. The only exceptions are non-equational
material generated by a document package, such as a tensor-network diagram or
table, and a short formula retained in inline math because no display is
warranted. These exceptions do not permit an unlabelled displayed equation.

Treat formulas as parts of sentences and punctuate them accordingly. Short
definitions may remain inline when that placement makes the sentence clearer.

\section{Formal traceability}

Source citations and external references belong in running prose, whereas
\Lean{} declarations remain supporting footnotes and never substitute for
source statements or citations. Write formal identifiers with \verb|\leanid| and
place them in footnotes unless the formal declaration itself is the mathematical
object under discussion. For example, the prose may say that the formal theorem
is the same-dimensional gauge-phase component of a source lemma, while a
footnote identifies the declaration.\footnote{Formal declaration:
    \leanid{MPSTensor.gaugePhaseEquiv_of_overlap_norm_tendsto_one_of_irreducible_TP}
    in \doclink{TNLean/MPS/SharedInfra/GaugePhase}.}
Do not put formal identifiers in displayed equations merely as traceability
data.

Explain a replacement theorem in ordinary mathematics before identifying its
formal declaration. Verify its hypotheses in the notation of the note. A
formal identifier or source-file link does not replace either step.

\section{Shared commands}

All notes in this directory must load
\path{docs/paper-gaps/command.tex}. Shared notation belongs there. In
particular, use the shared commands \verb|\tr| and \verb|\leanid| to produce
$\tr(X)$ and \path{someDeclaration}; do not redefine them in an individual
note. This rule keeps the documents visually consistent and prevents
contradictory local conventions.

\section{Naming}

A note's file name has the form \path{<key>_<topic>.tex}. The key names the
source the note examines: the authors' initials followed by a two-digit year
(\path{cpsv16}, \path{pgvwc07}, \path{spwc10}), or an established short name
for the source---a book or review (\path{wolf}, \path{rmp}) or a paper known by
its subject (\path{mpu} for matrix product unitaries, \path{peps} for normal
PEPS). A key may also name a small fixed collection of companion papers
examined together (\path{cpsv17}). A note that audits the formal development
itself, with no single external source, takes the key \path{tnlean}. The
registry of keys, with the source named by each key, is maintained in the
repository configuration \path{texra-blueprint.toml}; the index generator
rejects an unregistered key.

The topic part states the mathematical subject, not the history of the note. A
file name is permanent once another document cites it. Version suffixes are not
used because a note is revised in place and its history is preserved by the
repository.

Blueprint prose cites a note like any other source, with
\verb|\cite{gap:<slug>}|; the bibliography entry carries the published URL of
the note, whether the note lives in this repository or another one. Lean
docstrings and comments use the repository path
\path{docs/paper-gaps/<slug>.tex}.

\section{Classification}

The conclusion of a note must give a verdict. The available classifications
are a notational clarification, a local correction, a scope restriction, an
unfaithful theorem, a source claim that is false as printed, and an open gap.
The classification follows the mathematics, not the amount of formal work
required.

Each note declares its classification and status machine-readably with
\verb|\gapnote{<kind>}{<status>}| directly after \verb|\maketitle|. The kind
is one of \path{clarification}, \path{local-correction},
\path{scope-restriction}, \path{unfaithful}, \path{false-source}, and
\path{open-gap}. The status is one of \path{open}, \path{wip},
\path{resolved}, and \path{historical}. The status \path{wip} marks a gap whose
elimination is actively underway and still counts as live. Severity is derived
from the kind and is never stated separately. A resolved note keeps its
classification and changes only its status; the published index lists live
notes first and dims resolved notes.

A local correction records a small adjustment whose corrected statement is
still available in the cited argument. A scope restriction records a theorem
proved only in a special case of the source theorem. An unfaithful theorem
relies on a hypothesis or intermediate result that the cited source does not
provide. Such a theorem must carry an in-source marker pointing to the
paper-gap note. A source claim is false as printed when a counterexample
satisfies the stated source hypotheses and contradicts its conclusion. This
classification concerns the cited assertion itself rather than a mismatch
introduced by the formal statement or proof.

\section{Revision}

A note may begin provisionally, but revisions must keep it as a coherent account
of the present mathematical understanding rather than a chronological log. If
later work proves that $\mathcal K$ follows from $\mathcal H$, rewrite the note
to show that derivation. If later work confirms that the extra hypothesis is
unavoidable, state the restricted theorem explicitly. Preserve stable equation
labels when their mathematical content is unchanged, so references from other
documents remain valid.

\bibliographystyle{alpha}
\bibliography{references}

\end{document}
